\documentclass[10pt]{article}
\usepackage[margin=20mm]{geometry}
\usepackage{amsmath,amssymb,fontspec,unicode-math,xurl,hyperref}
\IfFontExistsTF{DejaVu Serif}{\setmainfont{DejaVu Serif}}{\setmainfont{Latin Modern Roman}}
\setmathfont{latinmodern-math.otf}
\setlength{\parindent}{0pt}\setlength{\parskip}{6pt}\sloppy
\begin{document}
Author: AI-assisted scientific verification; human decision pending\par
Date: 24 October 2026\par
Title: 0065 | Scientific Verification Report\par
I25-40174 | Gressmann, Chen, Xie, Amato and Rauchwerger | Position: It Is Time We Test Neural Computation In Vitro | PMLR 267:81393-81409 (2025). Reviewed publisher PDF: 17 pages.\par
Overall finding: Minor. The MNIST presentation-time estimate needs a corrected value or an explicit subset size. No substantial core defect was demonstrated in the inspected biological examples or proposed control framework. Human scientific decision pending; author adoption not requested.\par
\# 1. Source identity, version and audit scope\par
The original papers/0065/0065.pdf was hashed and matched to the controller acquisition receipt before review. The title, author list and proceedings identity agree with PMLR. A current publisher-linked PDF was independently downloaded and found byte-identical to the reviewed 4,090,192-byte, 17-page original. Original bytes were preserved.\par
{\ttfamily\footnotesize SHA-256: 0ee7f5189d626f915cef175b221c94f63c1af68148fa6429ed5c16b27fc1c4f7}\par
All main-text sections (pp.1-9) were inspected; references on pp.10-17 were checked selectively. The paper has no separate experimental appendix. The audit focuses on the biological proof-of-concept evidence, the proposed control/optimization formulation, the distinction between control and mechanistic identification, and the explicit numerical timing example. Figures 1 and 3 and the p.7 timing passage were visually checked.\par
Claim group | Verification and result\par
Biological examples and accessible platforms; Sections1-3 | Check published Kagan2022, Jordan2024, Zhang2024 and Elliott2023 sources, including cell provenance and stated limits. No demonstrated core contradiction.\par
Control model and optimization; Sections3.1-3.3 | Inspect equations, Figure1 and qualifications. Record Markov-state, computation-attribution and interface requirements as boundaries of the proposal.\par
MNIST timing; Section3.2,p7 | Check the actual training-label file header and exact arithmetic. Minor unqualified estimate / missing subset size.\par
Alternative views and long-term implications; Sections4-6 | Distinguish disclosed technical/ethical limitations and energy-efficiency expectations from demonstrated results. No new experiment or proven convergence result is claimed.\par
The PMLR and author-lab records did not expose an exact benchmark implementation for this position paper. No matching arXiv version was identified in the inspected records. Direct OpenReview access encountered a browser-verification wall, so a full OpenReview revision comparison was not performed. These are bounded search results, not assertions that no related material exists anywhere.\par
\# 2. Problem 1 (Minor): the MNIST epoch duration requires correction or a stated subset\par
Location and claim. In Section 3.2 on p.7, the paper assumes 500 milliseconds of presentation per image and estimates almost seven hours for one MNIST training epoch. No 50,000-example subset or train-validation partition is stated in the inspected main text.\par
Dataset check. The standard training split contains 60,000 examples. The audit independently downloaded the original-format train-labels IDX gzip file from the Google-hosted dataset mirror, decompressed it, and read its big-endian header: magic number 2049 and item count 60,000. The decoded file has exactly 60,008 bytes, consistent with an eight-byte header plus one label per example. Official Keras dataset documentation independently specifies the same training size.\par
\[T=N\times 0.5\,\mathrm{s},\qquad
T_{60{,}000}=30{,}000\,\mathrm{s}
=\frac{25}{3}\,\mathrm{h}=8\,\mathrm{h}\,20\,\mathrm{min}.\]\par
Assumed presented examples | Presentation time at500ms each\par
60,000: complete standard training split | 30,000s = 8h20m00s\par
50,000: possible training subset | 25,000s = 6h56m40s\par
Strongest alternative reading. Reserving 10,000 of the training examples for validation would yield a 50,000-example training subset and explain almost seven hours. That interpretation is possible but undocumented here. The evidence establishes a mismatch with the standard full-training-split reading, not an error under every imaginable partition.\par
Correction. State 8 hours 20 minutes for the full 60,000-example split, or explicitly state a 50,000-example training subset if that is intended. Specify that this is presentation time only, before reset intervals, pauses, data processing, repeated trials or other overhead.\par
Impact and severity. This is a bounded numerical/setup-description issue. Correcting the full-split estimate increases the required time and preserves the paper’s point that ordinary digital benchmarks can consume substantial biological recording time. It does not invalidate the proposed research program. Severity is Minor because no material central-claim damage has been demonstrated.\par
Replay. verification/0065/check\_mnist\_timing.py reproduces the download, header check and exact arithmetic; mnist\_timing\_check.json and mnist\_source\_manifest.json record results and hashes. No MNIST training run or wet-lab experiment was performed. The source gzip SHA-256 is `3552534a``0a558bbe``d6aed32b``30c495cc``a23d567e``c52cac8b``e1a0730e``8010255c`.\par
\# 3. Biological evidence: supported claims and their limits\par
The p.1 footnote narrows the paper’s use of in vitro neurons to iPSC-derived neurons. This makes source-cell provenance relevant, but does not automatically disqualify every citation involving other biological preparations, particularly where the cited contribution is instrumentation or a comparison with in vivo control.\par
Primary source and target use | What was verified; limit\par
Kagan etal.2022,DishBrain; targetSection3,p5 | The published study includes human iPSC-derived neurons as well as primary mouse cortical neurons, with closed-loop Pong-like gameplay and control conditions. This supports the target’s basic-gameplay proof of concept. It does not establish persistent across-day learning or a uniquely identified biological learning rule.\par
Jordan etal.2024,Neuroplatform; targetpp2,8 andFigure1 | The methods specify neural stem cells derived from human iPSCs, and the paper documents remote electrophysiological access. This supports the availability example. The platform source also reports capacity and experimental limitations; universal access, price and task capability were not audited here.\par
Zhang etal.2024,MindInVitroplatforms; targetpp2,7 | The demonstrations include mouse embryonic-stem-cell-derived and primary preparations, rather than establishing all cited results in human iPSC neurons. The target cites interface availability and recording practicalities. Transfer to its narrower cell scope remains something to validate; the source difference alone does not refute the instrumentation rationale.\par
Elliott etal.2023,remote closed-loop infrastructure; targetp2 | The published experiments use mouse embryonic stem-cell-derived cultures. They provide a remote-infrastructure precedent; Jordan separately supports human-iPSC platform availability. Neither establishes mature task computation across arbitrary preparations.\par
The narrow DishBrain claim in the target does not inherit every stronger interpretation suggested by the cited article’s title. In particular, the review does not treat basic task adaptation as established phenomenal consciousness, general intelligence, durable memory or a proof of a particular learning mechanism. The target’s later discussion acknowledges ethical uncertainty and technological immaturity.\par
Li et al. (2024) is identified in the target as an in vivo C. elegans example, not an iPSC culture experiment. The publisher abstract supports optogenetic integration of a deep-RL agent with the animal on a target-finding task. Full subscription-only methods and the exact offline Soft Actor-Critic training details were not independently reproduced in this audit.\par
Primary published XML snapshots and extracted text for the four platform/culture papers are preserved under sources/0065/bio\_primary/. The check examined what the papers report; it did not reproduce original spike recordings, culture-level statistics, biological replicates, plasticity assays or experimental controls.\par
\# 4. Control, computation and mechanistic inference: verification boundaries\par
Section 3.1 models the observed neural system with a transition law and a stimulation policy. Such notation is a useful starting point. If two histories produce the same current observation but different future responses, however, the current observation alone is not a sufficient Markov state. A practical formulation may need history, a learned latent state or a belief state. The paper’s proposed sequence models allow that interpretation, so this is a specification requirement rather than a counterexample to a universal theorem supplied in the paper.\par
Task success also does not by itself determine where the computation occurred. A simple logical example is a digital encoder that computes an XOR label, a fixed biological channel that relays its input, and an identity decoder. The combined system predicts XOR perfectly without the biological channel learning XOR. A four-row exact truth table is saved in model\_exact\_boundaries.json. This illustrates a measurement boundary, not a biological simulation or a falsification of the proposal.\par
Counterevidence is explicit. Section 3.1 permits reservoir/random-projection uses of the culture; Section 3.2 calls for objectives that incentivize nontrivial biological use. Thus the manuscript does not prove that any high task score establishes biological computation, internal plasticity or mechanistic identification. Future implementations would need suitable encoder/decoder controls, biological perturbations and attribution criteria to make those stronger inferences.\par
The suggestion that the artificial policy should ideally approach an identity map is an aspiration. Its mathematical meaning requires matching interfaces or defined encodings between task inputs, recorded activity and stimulation channels; the earlier map from state space to action space does not by itself define that identity. No convergence theorem, limit proof or trained implementation is supplied, and none was assumed in this audit.\par
The control-to-understanding argument is appropriately narrower when read as discovery of exploitable response properties. It does not identify a unique internal mechanism from arbitrary control success. The report preserves this limited reading rather than claiming that the paper has established or failed a stronger identifiability theorem.\par
Simulation fidelity, transfer and automatic differentiation are presented as research opportunities. The manuscript acknowledges simulation-to-experiment gaps and requires eventual testing on actual cultures. Results from a differentiable simulator or formal spiking model are not treated here as measured gradients through living tissue.\par
\# 5. Overall assessment, unperformed work and provenance\par
The inspected evidence supports a research agenda for experimenting with living neural systems through learned controllers. Cost, maturity and longer-term energy benefits remain conditional expectations. This review does not certify net energy savings for a complete installation including stimulation, recording, digital controllers and culture-maintenance infrastructure.\par
Not performed: wet-lab replication; reanalysis of the cited biological experiments or participant/culture-level statistical data; training a controller; executing author code; validating an exact benchmark, learning rule or causal mechanism; formal convergence/identifiability proofs; exhaustive reference or historical-timeline checking; and a current legal or ethical-policy audit. No general repair or implemented replacement system is claimed.\par
Independent AI checks of the biology and modeling portions, source hashes, acquisition outcomes and reviewer-written exact checks are saved in verification/0065/ and sources/0065/. The model examples are labeled boundary demonstrations, not newly discovered failures of a published algorithm. Original-source preservation and the five required deliverables are checked in completion\_checkpoint.json.\par
The earlier report PDF was generated with Lumen HTML/CSS from report\_content.json. The present delivery is compiled from this self-contained TeX using the existing Tectonic runtime and a task-local offline cache. Current compilation and visual-QA results are recorded in acceptance.json and the revision provenance record. Human scientific review and final scientific decisions remain pending.\par
Reviewed publisher record: https://proceedings.mlr.press/v267/gressmann25a.html\par
MNIST training-split documentation: https://keras.io/api/datasets/mnist/\par
Kagan etal.2022,published primary study: https://doi.org/10.1016/j.neuron.2022.09.001\par
Jordan etal.2024,remote platform: https://doi.org/10.3389/frai.2024.1376042\par
Zhang etal.2024,instrumentation: https://doi.org/10.1002/advs.202306826\par
Elliott etal.2023,infrastructure: https://doi.org/10.1523/ENEURO.0308-23.2023\par
Li etal.2024,in vivo control example: https://doi.org/10.1038/s42256-024-00854-2\par
\end{document}
